\section{Dataset \& Preprocessing}

\subsection{Dataset overview}

The dataset is a Fruits-360-like image collection organized as an \code{ImageFolder}-compatible directory structure:
\begin{itemize}
\item \code{data_final/train/{apple,banana,orange}}
\item \code{data_final/val/{apple,banana,orange}}
\item \code{data_final/test/{apple,banana,orange}}
\end{itemize}

The raw images are 100×100 pixels (verified via system tooling), and the training pipeline resizes them to \textbf{64×64} to match the generator and classifier architectures.

\subsubsection{Class balance and split integrity}

All three classes are perfectly balanced within each split (Table~\ref{tab:dataset-split-sizes}), which simplifies interpretation: accuracy changes are less likely to be driven by class imbalance.

The code assumes that the split directories (\code{train/}, \code{val/}, \code{test/}) contain disjoint images. The directory-based split helps prevent leakage, but this report does not perform hash-based duplicate detection across splits; thus, strict non-overlap is an assumption.

\subsubsection{Split sizes}

Exact image counts (exported in \code{reports/tables/dataset_summary.csv}) are:

\begin{table}[H]
\centering
\caption{Dataset split sizes by class.}
\label{tab:dataset-split-sizes}
\begin{tabular}{@{}lrrrr@{}}
\toprule
Split & Apple & Banana & Orange & Total \\
\midrule
train & 1300 & 1300 & 1300 & 3900 \\
val & 159 & 159 & 159 & 477 \\
test & 492 & 492 & 492 & 1476 \\
\bottomrule
\end{tabular}
\end{table}

\subsection{Real vs synthetic examples}

To make the real--synthetic gap tangible, the report includes deterministic exemplar images.

\textbf{Real exemplars} (lexicographically first file per class from \code{data_final/train/...}):

\begin{figure}[H]
\centering
\begin{subfigure}[t]{0.31\linewidth}
\centering
\includegraphics[width=\linewidth]{real_apple.jpg}
\caption{Apple}
\end{subfigure}\hfill
\begin{subfigure}[t]{0.31\linewidth}
\centering
\includegraphics[width=\linewidth]{real_banana.jpg}
\caption{Banana}
\end{subfigure}\hfill
\begin{subfigure}[t]{0.31\linewidth}
\centering
\includegraphics[width=\linewidth]{real_orange.jpg}
\caption{Orange}
\end{subfigure}
\caption{Real exemplar images from the training split.}
\label{fig:real-exemplars}
\end{figure}

\textbf{Synthetic exemplars} (\code{*_synth_00000.png} from \code{data_synth/...}):

\begin{figure}[H]
\centering
\begin{subfigure}[t]{0.31\linewidth}
\centering
\includegraphics[width=\linewidth]{synth_apple.png}
\caption{Apple}
\end{subfigure}\hfill
\begin{subfigure}[t]{0.31\linewidth}
\centering
\includegraphics[width=\linewidth]{synth_banana.png}
\caption{Banana}
\end{subfigure}\hfill
\begin{subfigure}[t]{0.31\linewidth}
\centering
\includegraphics[width=\linewidth]{synth_orange.png}
\caption{Orange}
\end{subfigure}
\caption{Synthetic exemplar images from the generated data pool.}
\label{fig:synthetic-exemplars}
\end{figure}

These single examples are not sufficient for evaluation, but they help qualitatively interpret the quantitative results.

\subsection{Preprocessing and augmentation}

\subsubsection{GAN training transforms}

GAN training uses:
\begin{itemize}
\item Resize to 64×64
\item Random horizontal flip (training only)
\item Tensor conversion + normalization to \([-1, 1]\) using mean=0.5, std=0.5 per channel
\end{itemize}

FID evaluation uses the same resize and normalization but disables augmentation.

\subsubsection{Classifier training transforms}

The classifier is trained with stronger augmentation:
\begin{itemize}
\item Resize to 64×64
\item Random horizontal flip
\item Random rotation up to ±15°
\item Color jitter (brightness and contrast)
\item Tensor conversion + normalization to \([-1, 1]\)
\end{itemize}

Test-time evaluation uses only resize + normalization.
